\section{Introduction}
\label{sec:intro}

Alignment research increasingly turns on null results. A red-team attack fails,
a safety property holds under stress testing, or a training change leaves a
behavioural metric unmoved. Researchers routinely read such a no-effect finding
as good news: the model is robust. But a null is one of the most ambiguous
outcomes an experiment can produce. ``Failure to reject $H_0$'' is, by
construction, \emph{not} evidence for $H_0$~\citep{dale2024confirming}, and the
same flat result can arise from at least three very different causes.

\para{Three causes of a null.} A null can mean (i) \emph{genuine robustness}---the
model truly is invariant to the manipulation; (ii) a \emph{weak intervention}---the
attack was too weak, the sample too small, or the metric too coarse to detect a
real effect; or (iii) a \emph{design flaw}---a confound that cancels a real
effect, evaluation awareness, a broken measurement, or a hypothesis whose logical
structure makes it unconfirmable in the first place~\citep{dale2024confirming,
santosgrueiro2026limits}. These three readings carry opposite implications for
what a researcher should do next, yet the standard statistical apparatus reports
the same $p > 0.05$ for all of them. Mis-reading a null wastes research effort
and, worse, can ship false confidence in a model's safety---declaring ``robust''
when the attack was simply too weak, or when a confound masked a real
vulnerability.

\para{Why existing tools are not enough.} Two literatures each solve half of the
problem. Bayesian model criticism can \emph{quantify} evidence for a null and
gauge whether an experiment had the power to detect an effect, through Bayes
factors, relative belief ratios, and posterior predictive
checks~\citep{rouder2009bayesian, moran2022ppn, robert2016demise}. But it is
blind to the logical structure of the experiment: a confound that cancels a real
effect looks like a clean null, and a two-sided point null is statistically
testable yet logically \emph{unconfirmable}~\citep{dale2024confirming}.
Conversely, symbolic and neurosymbolic reasoning can encode \emph{why} a null is
or is not interpretable---through confirmability topology and
identifiability/equivalence-class arguments~\citep{dale2024confirming,
santosgrueiro2026limits, tran2025neurosymbolic}---but cannot, on its own, tell a
genuinely-null effect from a real-but-underpowered one, because that distinction
needs sample size and measurement noise. And critically, \emph{no public dataset
labels the cause of a null}: every alignment-evaluation set diagnoses models, not
the epistemic status of a no-effect finding. The fusion of the two views, and a
benchmark to measure it, are both missing.

\para{Our approach.} We propose a hybrid Bayesian--symbolic diagnostic that
classifies the cause of a null result into \GR, \WI, or \DF, each with an
interpretable reasoning chain. A symbolic layer first gates confirmability and
identifiability, catching design flaws that statistical evidence cannot see.
Within the admissible structure, the framework requires \emph{all three}
interpretable adequacy checks---adequate sample, precise measurement, and
substantive intervention---to declare robustness; any failure routes the null to
weak intervention while naming the failed check. Because the two layers observe
disjoint, complementary inputs (sample size and noise on the Bayesian side;
structural flags and intervention strength on the symbolic side), only their
fusion holds all the signals needed to separate the three causes. We evaluate on
a literature-grounded, ground-truth-labelled generator, the only way to directly
measure three-way cause accuracy given the dataset gap above.

\para{Results preview.} On 3{,}000 held-out scenarios, the hybrid reaches a
three-way macro-F1 of $0.991$, against $0.871$ for symbolic-only, $0.489$ for
Bayesian-only, $0.326$ for equivalence-test-only, and $0.244$ for naive NHST---a
gain of $+0.12$ over the strongest single modality and more than $4\times$ over
the naive baseline, with non-overlapping bootstrap confidence intervals and
Holm-corrected McNemar tests at $p \ll 0.001$. The gain is mechanistic: the
Bayesian view never recovers a single design flaw ($\textsc{F1}=0.00$), while the
symbolic view mislabels $371/989$ underpowered cases as robustness; the hybrid
resolves both, leaving only $28/3{,}000$ errors. A frontier large language model
(GPT-4.1) given identical information scores $0.888$, competitive with the
symbolic baseline but well below the hybrid on the same items.

\para{Contributions.} We make the following contributions:
\begin{itemize}[leftmargin=*,itemsep=0pt,topsep=0pt]
  \item We formalise the interpretation of a null result in alignment experiments
        as a three-way cause-classification problem (\GR vs.\ \WI vs.\ \DF), and
        argue that resolving it requires \emph{both} statistical-adequacy signals
        and structural/identifiability signals.
  \item We propose a hybrid Bayesian--symbolic diagnostic that gates
        confirmability symbolically and quantifies evidence and power with a JZS
        Bayes factor, fusing the two into a calibrated, interpretable three-way
        classifier.
  \item We build a literature-grounded, ground-truth-labelled benchmark of null
        causes---to our knowledge the first---and show the hybrid beats
        Bayesian-only, symbolic-only, and two statistics-only baselines by wide,
        statistically significant margins, with a confusion analysis that pins
        each modality's blind spot.
  \item We add an external check with a frontier LLM expert judge (GPT-4.1) and a
        prior-sensitivity sweep, showing the structured framework adds value a
        strong general reasoner does not supply for free, and is robust to prior
        specification.
\end{itemize}

The rest of the paper reviews related work (\secref{sec:related}), describes the
framework, benchmark, and evaluation protocol (\secref{sec:method}), reports
results (\secref{sec:results}), discusses interpretation and limitations
(\secref{sec:discussion}), and concludes (\secref{sec:conclusion}).
